\section*{الفصل \ref{ch:b2:global-change} --- التغير العالمي والمحيط الحيوي}

\begin{solution}{exo:b2:global-change:1}
$\Delta T = \lambda C_{\text{المتراكم}}$ عند $\lambda = \qty{1.65}{\celsius}$
لكل \qty{1000}{GtC}: أي \qty{1.2}{\celsius} و\qty{1.65}{\celsius}
و\qty{3.3}{\celsius}.
\end{solution}

\begin{solution}{exo:b2:global-change:2}
درجات-الأيام: مجموع فائض الحرارة اليومية المتوسطة
فوق أساس على مدى الأيام، وهو \omterm{thm:b2:global-change:phenology}{الزمن الحراري} الذي يحكم
نماء النبات وذوات الحرارة المتغيرة. وعدم التزامن: حين يوقّت نوعان
متفاعلان فصليهما بإشارتين مختلفتين ويحرّك
الاحترار أحدهما دون الآخر --- فالخاطف، المستدل بطول
النهار في إفريقيا، يصل في تاريخه القديم فيجد ذروة
اليرقات، المستدلة بالحرارة، قد مضت منذ أسبوعين.
\end{solution}

\begin{solution}{exo:b2:global-change:3}
نحو \qty{150}{m} صعودًا و\qty{150}{km} نحو القطب لكل درجة. و
النوع الذي عند القمة لا أرض أعلى له: فخط حرارته يرتفع
فوق الجبل ويزول موئله، في حين تتقلص مساحة
كل نطاق مع الارتفاع في الطريق صعودًا.
\end{solution}

\begin{solution}{exo:b2:global-change:4}
\ensuremath{\mathrm{CO_2}} المنحل يكوّن حمض الكربون، وهو
يطلق شوارد هدروجين، فتهبط الحموضة بنسبة
لوغاريتم ضغط \ensuremath{\mathrm{CO_2}}. وشوارد
الهدروجين الزائدة تحوّل الكربونات إلى بيكربونات، فيهبط تركيز
الكربونات وتهبط حالة إشباع كربونات الكالسيوم.
والكائنات التي تبني الأراغونيت --- المرجان والجناحيات الأقدام وكثير من اليرقات
--- والكلسيت --- الطحالب الكلسية والمنخربات والرخويات ---
عليها أن تنفق أكثر لبناء أصدافها، وتفقدها دون الإشباع.
\end{solution}

\begin{solution}{exo:b2:global-change:5}
\qty{1.5}{\celsius}: $1.5/1.65\times 1000 = \qty{909}{GtC}$ إجمالًا،
و\qty{209}{GtC} باقية، أي 21 سنة. و\qty{2}{\celsius}:
\qty{1212}{GtC}، و\qty{512}{GtC} باقية، أي 51 سنة.
\end{solution}

\begin{solution}{exo:b2:global-change:6}
$\sqrt{2\times 150/0.12} = 50$ يومًا؛ والتقدم $1.5/0.12 = 12.5$
يوم.
\end{solution}

\begin{solution}{exo:b2:global-change:7}
\qty{3}{\celsius}: ترتفع خطوط الحرارة $3000/6.5 = \qty{460}{m}$؛ و
ينتقل النطاق إلى \qtyrange{1260}{1860}{m}، لكن الجبل ينتهي عند
\qty{1800}{m}: فتضيع \qty{60}{m} العليا من النطاق ويجلس الباقي
حيث تكون المساحة $2^{-460/300} = 0.35$ مما كانت --- أي نحو
ثلثي المساحة يزول. وعند \qty{5}{\celsius}: \qty{770}{m}؛ فيكون
النطاق \qtyrange{1570}{2170}{m}؛ ولا يوجد إلا \qtyrange{1570}{1800}{m}،
أي ثلث عرض النطاق عند $2^{-770/300} = 0.17$ من
المساحة --- أي دون عُشر المساحة الأصلية: فالنوع
زائل فعليًا.
\end{solution}

\begin{solution}{exo:b2:global-change:8}
عند 560: الحموضة $7.93$، وشوارد الهدروجين $\times 1.9$؛ وعند 800: $7.79$، و$\times
2.6$؛ وعند 1000: $7.70$، و$\times 3.1$.
\end{solution}

\begin{solution}{exo:b2:global-change:9}
$1 - (1 - h)^{0.25}$: أي \qty{26}{\%} و\qty{53}{\%} و\qty{82}{\%}. و
الفقد الفوري أصغر لأن البقية ما تزال تحمل
جمهرات معظم الأنواع؛ وهي أصغر من أن تدوم
فتنقرض على مدى عقود --- وهو \omterm{thm:b2:global-change:sar}{دين الانقراض}.
\end{solution}

\begin{solution}{exo:b2:global-change:10}
$c = 2\sqrt{0.7\times 20} = \qty{7.5}{km/yr}$؛ و\qty{1000}{km} في
نحو 130 سنة. وتنصيف $r$ أو $D$ يقسم $c$ على
$\sqrt 2$، إلى \qty{5.3}{km/yr}؛ فلا ينصّفها أي منهما --- ولتنصيف
السرعة يجب تربيع الجداء $rD$ نزولًا إلى الربع.
\end{solution}

\begin{solution}{exo:b2:global-change:11}
المعدل الخلفي: $8\times 10^{6}\times 10^{-6} = 8$ في السنة، أي 800 في
القرن؛ وعند ألف ضعف، 8000 في السنة، أي \num{800000} في القرن.
وعشرة آلاف موثقة في قرن هي نحو عشرة أمثال
المعدل الخلفي، لا ألف --- لكن التوثيق لا يغطي إلا
بضعة في المئة من الأنواع الموصوفة والمرصودة (الطيور
والثدييات وقليل من النبات)، والمعدل بينها هو فعلًا مئة
ضعف المعدل الخلفي أو أكثر؛ أما الأغلبية غير الموصوفة فمعظم
انقراضاتها لا يُرى، وتقدير الأنواع والمساحة هو
المقبض الوحيد. والحقيقة بين المعدود
والمقدَّر، ولذلك كان المجال المذكور بهذا الاتساع.
\end{solution}

\begin{solution}{exo:b2:global-change:12}
الحرارة تتناسب مع \omterm{prop:b2:global-change:tcre}{الانطلاقات المتراكمة}؛ فحين
تكون الانطلاقات الصافية صفرًا يكف المجموع المتراكم عن النمو ويكف
الاحترار كذلك. وفي \omterm{thm:b2:biogeochemical-cycles:box}{نموذج الصندوق}، عند انطلاقات صفر، يبدأ
الفائض الجوي بالهبوط إذ تمضي المصارف في عملها، وذلك
كان سيبرّد --- لكن المحيط، وهو ما يزال يدفأ نحو التوازن
مع \ensuremath{\mathrm{CO_2}} المرتفع، يمضي في تدفئة
السطح، ويكاد الأثران يتلاشيان: فتبقى الحرارة
حيث هي قرونًا. أما الناقص من الحجة: فغازات الدفيئة
الأخرى (فقصر عمر الميثان يعني أن احتراره
يهبط سريعًا متى توقفت الانطلاقات)، والهباءات التي تقنّع الآن
شيئًا من الاحترار وتزول في غضون أسابيع من الانطلاقات التي تصنعها،
والتغذيات الراجعة --- كربون \omterm{def:b2:living-soil:soil}{التربة} الصقيعية، وموت الغابات ---
التي قد تضيف انطلاقات خاصة بها.
\end{solution}

\begin{solution}{pb:b2:global-change:1}
\textbf{1.} المسار A: $700 + 80\times 10 = \qty{1500}{GtC}$. والمسار B: $700 +
\tfrac12\times 50\times 10 = \qty{950}{GtC}$.
\textbf{2.} المسار A: $1.65\times 1.5 = \qty{2.5}{\celsius}$. والمسار B:
$1.65\times 0.95 = \qty{1.6}{\celsius}$.
\textbf{3.} سنة 2050. المسار A: \qty{1000}{GtC}، و\qty{1.65}{\celsius}. والمسار B:
الانطلاقات $10\int_0^{30}(1 - t/50)\,\mathrm{d}t = 10(30 - 9) =
\qty{210}{GtC}$، فالمجموع 910، و\qty{1.5}{\celsius}.
\textbf{4.} $\lambda\times 100 = \qty{0.165}{\celsius}$ في العقد،
أي دون القيمة الملحوظة $0.2$ قليلًا (وهي تشمل الغازات الأخرى
وزوال تقنيع الهباءات).
\textbf{5.} نعم: \omterm{prop:b2:global-change:tcre}{فالانطلاقات المتراكمة} تكف عن النمو سنة 2070، ويكف
الاحترار كذلك، فيبقى عند \qty{1.6}{\celsius} ---
فالتناسب يعني أن الحرارة يقررها المجموع لا
المعدل.
\textbf{6.} $2/1.65\times 1000 = \qty{1212}{GtC}$؛ ويبلغه المسار A
بعد $512/10 = 51$ سنة، أي سنة 2071.
\textbf{7.} $\sqrt{2\times 200/0.1} = 63$ يومًا.
\textbf{8.} الاحترار نسبةً إلى 2020: المسار A $2.5 - 1.2 =
\qty{1.3}{\celsius}$، والمسار B \qty{0.4}{\celsius}. والتقدم $\Delta T/b$:
13 يومًا على A، و4 أيام على B.
\textbf{9.} الفجوة: على A $20 + 13 = 33$ يومًا؛ وعلى B 24 يومًا.
\textbf{10.} اليرقات متاحة من اليوم 15 إلى اليوم 25 بعد
الإيراق. فعلى B: يصل الطائر في اليوم 24، داخل النافذة بالكاد.
وعلى A: في اليوم 33، أي بعد آخر يرقة بثمانية أيام.
\textbf{11.} يتقدم الطائر $3\times 1.3 = 4$ أيام: فالفجوة $33 - 4 =
29$ يومًا --- وهي ما تزال خارج النافذة.
\textbf{12.} يستطيع الطائر العيش عند الحرارة الجديدة؛ وما لا
يستطيعه هو إطعام فراخه حين يكون الغذاء قد زال، لأن إشارته
وإشارة فريسته تستجيبان للاحترار استجابتين مختلفتين.
\textbf{13.} على A: $1.3/6.5\times 1000 = \qty{200}{m}$ صعودًا و
$1.3/0.65\times 100 = \qty{200}{km}$ نحو القطب. وعلى B: \qty{57}{m} و
\qty{57}{km}.
\textbf{14.} في 80 سنة تتحرك الشجرة \qty{8}{km}: فهي لا تستطيع
تتبع \qty{200}{km} (على A) وتقصر عن \qty{57}{km} (على B) أيضًا ---
فالأشجار تتخلف، وستبقى أمداؤها خارج التوازن مع المناخ
قرونًا على أي من المسارين.
\textbf{15.} على A: يكون النطاق \qtyrange{1700}{2000}{m}: فهو يبلغ
القمة بالكاد، ولا شيء أبعد يذهب إليه. والجبل يضيق
مع الارتفاع، فلا تغطي الشريحة نفسها إلا $2^{-200/300} = 0.63$ ضعف
مساحتها السابقة --- وأي احترار إضافي لا أرض له لينقلها
إليها.
\textbf{16.} على B: \qty{57}{m} صعودًا، وعامل المساحة $2^{-57/300} = 0.88$: أي
فقد \qty{12}{\%}.
\textbf{17.} ثمانية عقود عند \qty{30}{km}: أي \qty{240}{km}، وهي أكثر من
\qty{200}{km} اللازمة على A: فالسمكة تلحق (وتخرج
من مصيدتها القديمة).
\textbf{18.} عند الحافة السفلى يلاقي النوع منافسات
وممرضات ومفترسات صاعدة من الأسفل لم يلاقها
قبلُ، وتتفوق فيزيولوجيتها على فيزيولوجيته المتكيفة مع الحرارة القديمة:
فالتراجع تنافسي بقدر ما هو
فيزيولوجي.
\textbf{19.} على A: $\mathrm{pH} = 8.2 - 0.9\log_{10}(600/280) = 7.90$؛
وشوارد الهدروجين $10^{0.30} = 2.0$ ضعف ما كانت سنة 1750. وعلى B: $8.04$؛ أي $1.44$
ضعف.
\textbf{20.} $\Omega = 4/(\text{النسبة})$: أي $2.0$ على A، و$2.8$ على B. وعند 2
يظل المرجان يتكلس لكن ببطء، وهياكله أضعف و
شعابه تتآكل أسرع مما تنمو؛ وعند $2.8$ يبني، وإن كانت
حرارة المسار B ما تزال تبيّضه.
\textbf{21.} على A: $400\times 0.4^{0.25} = 318$، أي 82 مفقودًا. وعلى B:
$400\times 0.8^{0.25} = 378$، أي 22 مفقودًا.
\textbf{22.} على A: $318\times(1 - 0.05\times 1.3) = 298$. وعلى B:
$378\times(1 - 0.05\times 0.4) = 371$.
\textbf{23.} على A: $1 - 298/400 = \qty{26}{\%}$. وعلى B: \qty{7}{\%}.
\textbf{24.} على المسارين: الإزالة؛ أي 82 في مقابل 20 نوعًا على A، و22
في مقابل 7 على B --- فحد الاحترار في النموذج متواضع؛ وفي الواقع
يتفاعل الاثنان، إذ لا تستطيع غابة مشدَّفة إزاحة مداها.
\textbf{25.} المسار A: \qty{2.5}{\celsius}، والحموضة $7.90$، وربع
الأنواع الشجرية مفقود. والمسار B: \qty{1.6}{\celsius}، والحموضة $8.04$، و\qty{7}{\%}
مفقودة.
\end{solution}
